\documentclass[12pt]{article}
\usepackage{amsmath,amsfonts,amscd,amsthm}
\usepackage{hyperref}
\hypersetup{
    colorlinks=true,
    linkcolor=black,
    citecolor=black,
    urlcolor=blue
    }
\usepackage{enumitem}
\usepackage{dsfont}
\usepackage{xcolor}
\usepackage{tikz}
\usepackage{multirow,booktabs}
\usepackage{graphics,subfig}
\usepackage{setspace}
\usepackage{changepage}
\usepackage{multicol}
\usepackage[retainorgcmds]{IEEEtrantools}
\usepackage{geometry}
\usepackage{dcolumn}
\usepackage{cleveref}
\newlength{\tabcont}
\setlength{\parindent}{0.0in}
\setlength{\parskip}{0.05in}
\usepackage{natbib}
\geometry{top=2cm,bottom=2cm,left=3cm,right=2cm, headsep=0in}
\parindent 0in
\parskip 12pt
\linespread{1.5}

% The theoremstyle for a theorem goes right above the \newtheorem command
\theoremstyle{definition}
\newtheorem{definition}{Definition}[section]

\newtheorem{theorem}{Theorem}[section]
\newtheorem{corollary}{Corollary}[section]
\newtheorem{remark}{Remark}[section]
\newtheorem{assumption}{Assumption}[section]
\renewcommand{\baselinestretch}{1.5} 

\title{Master Thesis Topic}
\author{Fabian Christian Schmidt}
\date{\today}

\begin{document}

\begin{titlepage}
    \begin{center}
        \vspace*{0.5cm}
            
        \huge
        \textbf{On Functional Garch Models and Their Application to Intraday Value-At-Risk}
            
        \vspace{1cm}
        \large
        Master Thesis Presented to the\\
        Department of Economics at the\\
        Rheinische Friedrich-Wilhelm-Universität Bonn
            
        \vspace{1.5cm}
            
        In Partial Fulfillment of the Requirements of the Degree of\\
        Master of Science (M.Sc.)
            
        \vspace{4cm}
            
        \large
        Supervisor: Prof. Dr. Dominik Liebl

        \vspace{1cm}
        Submitted in March 2026 by:\\
        Fabian Christian Schmidt\\
        Matriculation number: 50092796
            
    \end{center}
\end{titlepage}
\tableofcontents
\thispagestyle{empty}
\newpage
\setcounter{page}{1}
\section{Introduction}
\label{sec:Introduction}
%\begin{itemize}
%    \item Start with explanation of the history of interest in risk measures
%    \item transition to GARCH models and their usage when estimating risk $\longrightarrow$ name most important papers regarding GARCH models
%    \item transition to Functional Data Analysis and then Functional GARCH and its applications (literature)
%    \item explain structure and goal of the thesis
%\end{itemize}
The interest in quantitative financial risk measures has a long history, with the earliest ideas of quantitative risk measures dating back to the early $20^{th}$ century. The most notable early contribution was the dissertation by Louis Bachelier, which introduced Brownian motion to finance. Risk management, in general, dates back to the $17^{th}$ century, with merchants using rudimentary forms of protection against trade risks. The goal of risk management should always be to limit the risk of (high) losses and default while not establishing too high levels of security such that the business becomes unprofitable. One such measure that is the most popular for measuring market risk, the risk associated with losses on balance sheet and off-balance sheet items, is value-at-risk. The concept of value-at-risk was first mentioned in \citet{Leavens1945} before being introduced by \citet{Markowitz1952} and \citet{Roy1952} and became the preferred method for estimating market risk with the Basel II Accord in June 2004. Value-at-risk became so popular for estimating  One approach to estimate value-at-risk is to make use of so-called (general) auto-regressive conditional heteroskedasticity models.\\
The GARCH model, a generalization of the previously published ARCH model by \citet{Engle1982}, was introduced by \citet{Bollerslev1986}. (G)ARCH models differ fundamentally from standard time series models such as ARMA or VAR models. Instead of modeling the conditional mean of returns, they model the conditional variance by relating squared returns to their own lagged values. This approach is motivated by key stylized facts of financial time series. While returns typically exhibit little to no autocorrelation, squared or absolute returns display strong and persistent autocorrelation. In addition, volatility tends to cluster over time, with large fluctuations occurring in a recurrent but non-periodic manner \citep{mandelbrot1963variation,taylor2011asset}. (G)ARCH models provide an appropriate framework for capturing these stylized facts. By explicitly modeling the variance of financial assets the GARCH model allows the user to get an understanding of the underlying process related to risk of a financial asset which allows the user to forecast value-at-risk.\\
With the rise of high frequency data available both in finance and other fields, appropriate models, and theoretical frameworks are needed to make use of this data. Functional Data Analysis deals with data that manifests itself as curves. Since high frequency data in finance, e.g. return data, is highly correlated methods and theoretical frameworks from the field of Functional Data Analysis seem appropriate for this task. There exist various monographs for a comprehensive introduction into the subject. \citet{Ramsay2005} focuses on the practical application of Functional Data Analysis while \citet{hsing2015theoretical} deals with the theoretical foundations of Functional Data Analysis. \citet{kokoszka2017introduction} gives a good overview of both while also dedicating a few chapters to more specialized topics like Functional time series.\\
\citet{hormann2013functional} make use of methods from the field of Functional time series to extend the ARCH model to the functional context. \citet{Aue2017} then do the same with the GARCH model. \citet{RICE20201023} then apply this model to forecasting intra-day value-at-risk return curves. The latter two papers is what this thesis is about.\\
The aim of this thesis is to examine both the theoretical as well as finite sample properties of the Functional GARCH model in the context of using Functional GARCH models to forecast intraday value-at-risk. To do so, in \Cref{subsec:ScalarGARCH}, I will first give a selected excerpt of the main results and models of both univariate and multivariate GARCH models. Second, in \Cref{subsec:FGARCHModel}, I will present the Functional GARCH model and its theoretical properties. Third, in \Cref{subsec:FGARCHEstimation}, two possible methods for estimating the parameters of the Functional GARCH model are introduced and their consistency and asymptotic distributions are investigated. Fourth, in \Cref{subsec:ValueAtRisk}, the concept of value-at-risk and a brief overview of its history is given. Fifth, in \Cref{sec:Simulation}, the finite sample properties of the Functional GARCH model and implications on forecasting intraday value-at-risk are examined via the use of Monte Carlo simulations. Additionally, the Functional GARCH model is compared to two multivariate GARCH models presented in the preceding sections.
\section{Functional GARCH}
\label{sec:FGARCH}
\subsection{(Scalar) GARCH and Multivariate GARCH}
\label{subsec:ScalarGARCH}
Before I turn to the functional GARCH model, I will focus on the scalar version of the GARCH model, both in the univariate and multivariate cases.\\
An univariate GARCH(p,q) model can be described by its first two moments:
\begin{definition}[GARCH(p,q) Process]
    A process $(y_t)$ is called a GARCH(p,q) process if its first two conditional moments exist and satisfy:
    \begin{enumerate}
        \item $\mathrm{E}(y_t \,|\, y_u,u<t)=0, \quad \forall t\in\mathbb{Z}$.
        \item There exist constants $\delta,\alpha_i,i=1,...,q$ and $\beta_j,j=1,...,p$ such that\\
        \begin{align}
            \sigma_t^2=\mathrm{Var}(y_t\,|\,y_u,u<t)=\delta+\sum_{i=0}^q\alpha_iy_{t-i}^2+\sum_{j=1}^p\beta_j\sigma_{t-j}^2
        \end{align}
    \end{enumerate}
\end{definition}
A helpful representation of the GARCH(p,q) process can be obtained by defining an error term $\nu_t=y_t^2-\sigma_t^2$ and replacing $\sigma^2$ with $y^2-\nu$
\begin{align}
\label{ARMAstructure}
    y_t^2=\delta+\sum_{i=0}^r\overbrace{(\alpha_i+\beta_i)}^{\text{$\kappa_i$}}y_{t-i}^2+\nu_t-\sum_{j=1}^p\beta_j\nu_{t-j}=\delta+\sum_{i=0}^r\kappa_iy_{t-i}^2+\nu_t-\sum_{j=1}^p\beta_j\nu_{t-j}
\end{align}
where $r=\mathrm{max}(p,q)$ with $\alpha_i=0$ ($\beta_j=0$) if $i>q$ ($j>p$). From this representation, it can be seen that a GARCH(p,q) process follows an ARMA(r,p)-like structure for the squared values $y_t^2$. This insight also leads to the GARCH model being referred to as a generalization of the ARCH model which follows an AR-structure for the volatility equation. It can be shown, similarly to the AR and ARMA models, that a GARCH(p,q) process is equivalent to an ARCH($\infty$) process \citep{francq2019garch}. Note that $(\nu_t)$ is not i.i.d. and therefore cannot be used to examine the existence of a stationary solution of the GARCH(p,q) process.\\
A more restrictive yet often useful definition of a GARCH process is that of a strong GARCH(p,q) process. A strong univariate GARCH(p,q) process can be represented as follows:
\begin{gather}
    y_t = \sigma_t\varepsilon_t \\
    \sigma_t^2 = \delta + \sum_i^p \alpha_i y_{t-i}^2 + \sum_j^q \beta_j \sigma_{t-i}^2
    \label{Eq2StrongGARCH}
\end{gather}
where $(\varepsilon_t)$ is $\mathrm{i.i.d.}$ Additionally, it is often assumed that $E[\varepsilon_t|\varepsilon_j, j<t]=0$ and $\mathrm{Var}[\varepsilon_t]=1$.\\
I will focus on the strong GARCH(1,1) process since an equivalent functional version of this process is also used in \citet{Aue2017} as well as \citet{RICE20201023}.
The strong GARCH(1,1) process is the simplest specification of the strong GARCH(p,q) processes. Nevertheless, this specification is also one of the most widely used in practice.\\ % Add reference!!!
Before listing the conditions for a stationary solution to the strong GARCH(1,1) process to exist, note that the second equation of the strong GARCH(1,1) process can be rewritten as follows:
\begin{align}
    \sigma_t^2 = \delta + \gamma_{t-1} \sigma_{t-1}^2, \quad \text{where }\gamma_t=\alpha\varepsilon_t^2+\beta
\end{align}
From this representation, it can be seen that the volatility equation of the strong GARCH process can also be represented as an AR model with a random coefficient. This representation is useful for studying the existence of a stationary and ergodic solution.\\
It can be shown that there exists a strictly stationary, non-anticipative, ergodic solution to the strong GARCH(1,1) process if the following condition is satisfied:
\begin{theorem}[Strict stationarity of the strong GARCH(1,1) process]
\label{stationarityscalarGARCH11}
    If
    \begin{gather}
        -\infty\leq\mathbb{E}\,\ \mathrm{log}\{\overbrace{\alpha\varepsilon_t^2+\beta}^{\gamma_t}\}<0,
    \end{gather}
    then the infinite sum
    \begin{gather}
        h_t = \delta+\sum_{i=1}^{\infty}\Gamma_{t,i}\delta, \quad \text{where }\Gamma_{t,i}=\gamma_{t-1}...\gamma_{t-i}
    \end{gather}
    converges almost surely, and the process $(y_t)$ defined by $y_t=\sqrt{h_t}\varepsilon_t$ is the unique strictly stationary solution of model (2.7). This solution is non-anticipative and ergodic.
\end{theorem}
To prove \Cref{stationarityscalarGARCH11}, note that \labelcref{Eq2StrongGARCH} can be solved recursively to attain $h_t$. Showing that $h_t$ converges involves an application of the Cauchy Root Test as well as the strong law of large numbers. Uniqueness then follows by contradiction, assuming the existence of another solution and showing that it must coincide with the original solution.\\
Ergodicity describes the property of a stochastic process that time averages along a single, sufficiently long realization converge to the same values as ensemble averages taken over many realizations. A sufficient and particularly strong way to guarantee ergodicity is geometric ergodicity, which requires not only convergence to a unique stationary distribution but also convergence at an exponential rate. It can be shown that $(\sigma_t;t\in \mathbb{Z})$ is geometrically ergodic under mild assumptions meaning that initial conditions are forgotten exponentially fast and sample averages converge to their population counterparts. $(y_t;t\in\mathbb{Z})$ is geometrically $\beta$-mixing which implies (weak) ergodicity. To prove that the GARCH(1,1) process is geometrically ergodic first note that the volatility equation of the GARCH process can be rewritten into a Markov Chain representation. One can then make use of a Theorem from \citet{FeignTweedie1985}. They show that a Markov Chain that is also an aperiodic, irreducible Feller Chain is geometrically ergodic if it satisfies an energy condition. Proving that $(y_t;t\in\mathbb{Z})$ is geometrically $\beta$-mixing follows from the geometric ergodicity of $(\sigma_t;t\in\mathbb{Z})$ and the fact that $(\varepsilon_t;t\in\mathbb{Z})$ is i.i.d. $(y_t;t\in\mathbb{Z})$ is not geometrically ergodic since it is not Markovian. The complete proof can be found in the Appendix.\\ % Add a sentence on ergodicity and non-anticipativeness Maybe try using Cleveref here if there is time.
%The more general proof for strong GARCH(p,q) processes follows a similar structure
Since in practice, only $y_t$ is observed, the parameters $\delta$,$\alpha$, and $\beta$ have to be estimated. GARCH models are typically estimated using (quasi-)maximum likelihood estimation. (Quasi-)maximum likelihood estimation is preferred for GARCH models over least-squares estimation due to least-squares requiring higher moments and being less efficient than (quasi-)maximum likelihood estimation \citep{francq2019garch}.\\
The likelihood equations for the GARCH(1,1) process are as follows
\begin{gather}
    L_n(\theta;y_1,..,y_n)=\prod_{t=1}^n\frac{1}{\sqrt{2\pi \tilde{\sigma}_t^2}}\exp\left(-\frac{y_t^2}{2\tilde{\sigma}_t^2}\right); \quad \tilde{\sigma}_t^2=\tilde{\sigma}_t^2(\theta)=\delta+\alpha y_{t-1}^2+\beta\tilde{\sigma}_{t-1}^2.
\end{gather}
Here $\theta$ is a vector containing the estimates $\delta$, $\alpha$ and $\beta$ and $\tilde{\sigma}_t^2$ is calculated recursively. The true parameters are $\theta_0=(\delta_0,\alpha_0,\beta_0)^\top$.\\
Taking the logarithm, the likelihood optimization problem can be simplified to minimizing the following equation
\begin{gather}
    \mathcal{I}_n=\frac{1}{n}\sum_{t=1}^n \frac{y_t^2}{\tilde{\sigma}_t^2}+\log \tilde{\sigma}_t^2.
\end{gather}
%Taking the first derivative leads to the following moment condition
%\begin{gather}
   %\frac{1}{n} \sum_{t=1}^n \{y_t^2-\tilde{\sigma}_t^2\}\frac{1}{\tilde{\sigma}_t^4}\frac{\partial \tilde{\sigma}_t^2}{\partial \theta}=0.
%\end{gather}
From here, parameters can be found by using an optimization algorithm, e.g. the Newton-Raphson method.\\
It can be shown that the quasi-maximum likelihood estimator is consistent and asymptotically normal with the asymptotic variance of the estimator being equal to the inverse of the Fisher Information matrix when $(\varepsilon_t; t\in\mathbb{Z})$ is Gaussian and the inverse of the adjusted Fisher Information matrix when it is not. The details of the proof are omitted here for the sake of conciseness.\\
Next, I will give a short exposition of multivariate GARCH processes since, in practice, they are very closely related to Functional GARCH processes. Defining a multivariate GARCH process is not as straightforward as for the univariate case. This is the case since adding more processes to the multivariate GARCH process does not only require estimating another conditional variance equation, but also the conditional covariances between the processes. This quickly leads to estimation becoming statistically infeasible if estimation stays unconstrained. Therefore, the conditional covariance matrix $H_t$ follows an explicit specification. In the following, I will present two of these specifications since they are closely connected to the Functional GARCH model. Throughout this I will assume that $y_t,\varepsilon_t\in \mathbb{R}^m$ and that $\varepsilon_t$ is an i.i.d. vector.\\
First, I consider the Principal Components GARCH (PC-GARCH) model, which is defined by
\begin{gather}
\label{PCGARCHyequation}
    y_t = H_t^{\frac{1}{2}}\varepsilon_t,\\
    \label{PCGARCHHequation}
    H_t = P \Lambda_t P^{\top}.
\end{gather}
Here, $P$ is an $m \times m$ orthogonal matrix, while $\Lambda_t$ is an $m \times m$ diagonal matrix with diagonal elements $\text{diag}(\lambda_{1,t}, \lambda_{2,t}, \ldots, \lambda_{m,t})$. The matrix $P$ is typically estimated using Principal Component Analysis, and the diagonal elements of $\Lambda_t$ are assumed to follow univariate GARCH(1,1) processes:
\begin{gather}
\label{PCGARCHlambdaequation}
    \lambda_{i,t} = \delta_i + \alpha_i f_{i,t-1}^{2} + \beta_i \lambda_{i,t-1}, \quad i = 1, \ldots, m.
\end{gather}
In this specification, $f_{i,t}$ denotes the $i$-th component of the transformed series $f_t = P^\top y_t$. \labelcref{PCGARCHlambdaequation} is estimated using quasi-maximum likelihood estimation, although alternative specifications for $\lambda_{i,t}$ are also possible.\\
From the representation in \labelcref{PCGARCHHequation}, it follows directly that the PC-GARCH model decomposes the conditional covariance matrix $(H_t)$ into a set of constant eigenvectors and time-varying eigenvalues. As discussed in \Cref{sec:Simulation} and in the Appendix, the PC-GARCH model is closely related to the discretized Functional GARCH model, in which the underlying functions are observed on a finite grid.\\
The second multivariate GARCH specification considered is the Constant Conditional Correlation (CCC) GARCH model, which is defined as
\begin{gather}
    y_t = H_t^{\frac{1}{2}}\varepsilon_t,\\
    \label{CCCGARCHHequation}
    H_t = D_t R D_t^{\top}.
\end{gather}
Here, $D_t$ is an $m \times m$ diagonal matrix with diagonal elements $\text{diag}(\sqrt{h_{1,t}}, \sqrt{h_{2,t}}, \ldots, \sqrt{h_{m,t}})$, and $R$ is an $m \times m$ constant correlation matrix. The diagonal elements of $D_t$ are assumed to follow a GARCH(1,1)-type specification given by
\begin{gather}
\label{CCCGARCHsmallhequation}
    h_t = \delta + A (y_{t-1} \odot y_{t-1}) + B h_{t-1},
\end{gather}
where $\odot$ denotes the elementwise (Hadamard) product, $\delta$ is an $m$-dimensional vector, and $A$ and $B$ are $m \times m$ parameter matrices. \labelcref{CCCGARCHHequation} shows that $H_t$ consists of time-varying, variable-specific conditional variances combined with a constant correlation structure. Additionally, as can be seen from \labelcref{CCCGARCHsmallhequation} the variable specific conditional variances evolves dynamically, depending not only on each variable’s own past but also on the past of other variables. As discussed in \Cref{sec:FGARCH}, the Functional GARCH model can be interpreted as a semi-strong CCC-GARCH model with $R = \mathbb{I}$, where $\mathbb{I}$ denotes the identity matrix, after projection onto a finite set of basis functions.
\subsection{Model}
\label{subsec:FGARCHModel}
The next section will introduce the functional GARCH model and highlight its properties. I will start by defining the model itself.\\
A functional GARCH process is defined as follows:
\begin{definition}[Functional GARCH(p,q) Process]
\label{fGARCHpq}
A sequence of random functions $(y_i(t))$ is called a functional GARCH process if it satisfies the following equations:
    \begin{gather}
        y_i(t) = \sigma_i(t)\varepsilon_i(t)\\
        \sigma_i^2(t) = \delta(t)+\sum_{j=1}^p\alpha_jy_{i-j}^2(t)+\sum_{k=1}^q\beta_k\sigma_{i-k}^2(t)
    \end{gather}
\end{definition}
Here, $\delta(t)$ is a non-negative function, and $(\alpha_j)_{j\in (1,...,p)}$ and $(\beta_k)_{k \in (1,...,q)}$ are both integral operators $\alpha_jx(t)=\int \alpha_j(t,s)x(s)ds$ ($\beta_jx(t)=\int \beta_j(t,s)x(s)ds$) that map non-negative functions to non-negative functions. $x$ here denotes an arbitrary element in $L^2[0,1]$. By convention, I restrict the intraday time variable to $t\in[0,1]$. This is without loss of generality since data that do not follow this convention can easily be converted into this format. $\int_0^1$ is therefore shortened to $\int$. The second time variable $i$ denotes the trading day. $(\varepsilon_i(t))_{i\in\mathbb{Z}}$ are i.i.d. functions.\\
From \Cref{fGARCHpq}, it is easily seen that the Functional GARCH(p,q) process follows a similar functional structure to the strong (scalar) GARCH process. As a consequence, \Cref{fGARCHpq} can also be rewritten into a similar ARMA-like structure as in the scalar version in \labelcref{ARMAstructure}.\\
To stick closely to the paper from \citet{Aue2017}, as well as \citet{RICE20201023}, I will focus on the Functional GARCH(1,1) model. That reduces the equations of the Functional GARCH(p,q) model to:
\begin{gather}
    \label{FunctionalGARCH11}
    y_i(t) = \sigma_i(t)\varepsilon_i(t)\\
    \label{FunctionalGARCH11volatility}
    \sigma_i^2(t) = \delta(t)+\alpha y_{i-j}^2(t)+\beta\sigma_{i-k}^2(t)
\end{gather}
In the following, conditions for the existence of a stationary solution to the Functional GARCH(1,1) model are discussed. \citet{Aue2017} derive the conditions both for the space $L^2[0,1]$ as well as for the space of continuous functions $C[0,1]$. I will first focus on solutions in $L^2[0,1]$.\\
To use similar notation to the scalar case, rewrite the volatility equation into:
\begin{gather}
\label{FGARCH11volatilityalternative}
    \sigma_i^2(t) = \delta(t) + \gamma_{i-1} \sigma_{i-1}^2(t), \quad \text{where }\gamma_i(t,s)=\alpha\varepsilon_i^2(t)+\beta
\end{gather}
The necessary condition for a stationary solution then is:
\begin{theorem}[Strict stationarity of the functional GARCH(1,1) process]
\label{StationarityFGARCH11}
    Let $\delta(t)\in L^2[0,1]$ be a non-negative function and $\alpha\in L^2[0,1]^2$ and $\beta \in L^2[0,1]^2$ be integral operators that map non-negative functions to non-negative functions. Let $\{\varepsilon_i(t):i\in \mathbb{Z}\}$ be i.i.d. random functions and $\varepsilon_0\in L^{4}[0,1]$ almost surely.
    \begin{enumerate}[label=(\roman*)]
        \item If
            \begin{gather}
            \label{FGARCH11StationaryCondition}
                \infty\leq \mathbb{E}[\mathrm{log\left\| \gamma_0 \right\|_\mathcal{S}}]<0
            \end{gather}
            then the infinite sum
            \begin{gather}
            \label{FGARCH11StationarySolution}
                h_i(t)=\delta(t)+\sum_{k=1}^\infty \Gamma_{i,k}\delta(t)
            \end{gather}
            converges almost surely, and the process $(y_i(t))$ defined by $y_i(t)=\sqrt{h_i(t)}\varepsilon_i(t)$ is the unique strictly stationary solution of the Functional GARCH(1,1) model in \labelcref{FunctionalGARCH11} and \labelcref{FunctionalGARCH11volatility}. This solution is non-anticipative and ergodic.
        \item If there is $\nu>0$ such that 
        \begin{gather}
        \label{FGARCH11StationarityMomentsCondition}
            \mathbb{E}[\left\| \gamma_0 \right\|_\mathcal{S}^\nu ]<1
        \end{gather}
        then there exists a unique, strictly stationary, and nonanticipative solution in $L^2[0,1]$ to the Functional GARCH(1,1) model and $\mathbb{E}[\left\| \sigma_0^2 \right\|_2^\nu]<\infty$.
    \end{enumerate}
\end{theorem}
Here $\left\|\right\|_\mathcal{S}$ is the Hilbert-Schmidt norm defined as
\begin{gather}
    \left\| \gamma_0 \right\|_\mathcal{S} = \left( \int\int \gamma^2_0(t,s)dt ds \right)^{\frac{1}{2}}
\end{gather}
and $\Gamma_{i,k}$ is the following random integral operator
\begin{gather}
    (\Gamma_{i,k}x)(t)=\int ...\int \gamma_{i-1}(t,s_1)\gamma_{i-2}(s_1,s_2)...\gamma_{i-k}(s_{k-1},s_k)ds_k...ds_2ds_1.
\end{gather}
Due to the close resemblance between this theorem and \Cref{stationarityscalarGARCH11}, the proof of \Cref{StationarityFGARCH11} follows analogous arguments. First, $h_i(t)$ is obtained by recursively solving \labelcref{FGARCH11volatilityalternative}. Second, the convergence of $h_i(t)$ is established by applying the Cauchy--Schwarz inequality in conjunction with the strong law of large numbers. Finally, the uniqueness of the solution is shown by exploiting the recursive structure of \labelcref{FGARCH11volatilityalternative}. Assuming the existence of an alternative solution implies that both solutions converge to the same limit, which in turn forces them to coincide, thereby yielding a contradiction. The full proof of \Cref{StationarityFGARCH11} can be found in the Appendix.\\
%My previous version:
%Due to the strong similarity of this theorem to \ref{stationarityscalarGARCH(1,1)}, proving \ref{StationarityFGARCH(1,1)} follows a similar idea. Firstly, $h_i(t)$ is derived by recursively solving (\ref{FGARCH(1,1)volatilityalternative}). Secondly, to prove the convergence of $h_i(t)$ an application of the Cauchy-Schwarz inequality as well as the strong law of large numbers is needed. Lastly, to prove the uniqueness of the solution, the recursive structure of (\ref{FGARCH(1,1)volatilityalternative}) is used. Assuming that another solution exists results in the two solutions converging to the same limit and therefore both solutions coincide resulting in a contradiction.
%Both \citet{Cerovecki2019} and \citet{Kuehnert2020} extend the stationarity result in (\ref{FGARCH(1,1)StationaryCondition}). \citet{Cerovecki2019} provide sufficient conditions under which the Functional GARCH model is stationary for the $L^2$ space for any length $p,q$. These conditions are sufficient but not necessary, as stating necessary conditions is more complicated in the functional setting since norms are not equivalent in the infinite dimensional setting. To do so, they make use of the top Lyapunov exponent \citep{Kingman1973}. They show that when applying their result to the FGARCH(1,1) model a milder condition than (\ref{FGARCH(1,1)StationaryCondition}) can be established. However, to do so, they require a different assumption on $\varepsilon_0$. \citet{Kuehnert2020} further expands conditions for stationarity for the FGARCH model to $L^p$ where $p\in[1,\infty)$ and Banach spaces.\\ My version
Both \citet{Cerovecki2019} and \citet{Kuehnert2020} extend the stationarity result in \labelcref{FGARCH11StationaryCondition}. \citet{Cerovecki2019} derive sufficient conditions for stationarity of the functional GARCH model in $L^2$ for arbitrary orders $p$ and $q$. These conditions are not necessary, as establishing necessity is more difficult in the functional setting due to the non-equivalence of norms in infinite-dimensional spaces. Their approach relies on the top Lyapunov exponent \citep{Kingman1973}. When their general result is applied to the FGARCH(1,1) specification, it leads to a milder condition than \labelcref{FGARCH11StationaryCondition}. However, this comes at the cost of imposing a different assumption on $\varepsilon_0$. \citet{Kuehnert2020} further generalizes stationarity conditions to $L^p$ spaces with $p \in [1,\infty)$, and more broadly, to Banach spaces.\\
To prove the ergodicity of the above solution, the following corollary is useful.
\begin{corollary}
    \begin{enumerate}[label=(\roman*)]
        \item
        \label{FirstPointCorollary}
        There is a functional $g$ such that
        \begin{gather}
            \sigma_i^2(t)=g(\varepsilon_i(t),\varepsilon_{i-1}(t),...).
        \end{gather}
        As a consequence, $(\sigma_i^2(t);i\in\mathbb{Z})$ and $(y_i(t);i\in \mathbb{Z})$ are Bernoulli shifts. This follows directly from \labelcref{FGARCH11StationarySolution}.
        \item
        \label{SecondPointCorollary}
        \labelcref{FirstPointCorollary} implies that both $(\sigma_i^2(t); i\in \mathbb{Z)}$ and $(y_i(t);i\in\mathbb{Z})$ are ergodic.
    \end{enumerate}
\end{corollary}
(Weak) ergodicity of Bernoulli shifts is a well established result. (Weak) ergodicity is proven by...\\
Unlike in the scalar version of the GARCH model proving geometric ergodicity and $\beta$-mixing for the Functional GARCH model is very difficult. That is because for infi However, (weak) ergodicity is not sufficient for the derivation of important statistical results like the central limit theorem. That is because information on how fast dependence decays is needed to establish the convergence of   Since (weak) ergodicity is not sufficient for the application of the functional central limit theorem, \citet{Aue2017} derive the following result.
\begin{corollary}
    \begin{enumerate}[label=(\Roman*)]
        \item
        \label{FGARCH11CLTCondition}
        For $\ell\in \mathbb{N}$ and $j,i\in\mathbb{Z}$, let $\varepsilon_{i,j}^{(\ell)}(t)$ be independent copies of $\varepsilon_0(t)$ and define
        \begin{gather}
            \sigma_{i,\ell}^2(t) = g\left(\varepsilon_i(t),...,\varepsilon_{i-\ell}(t),\varepsilon_{i-\ell-1,i-\ell}^{(i)}(t),\varepsilon_{i-\ell-2,i-\ell}^{(i)}(t),... \right).
        \end{gather}
        If \labelcref{FGARCH11StationarityMomentsCondition} is satisfied, then there exists a constant $0<\rho<1$ such that
        \begin{gather}
            \mathbb{E}\left[\left\|\sigma_i^2(t)-\sigma_{i,\ell}(t) \right\|_2^\nu\right]=\mathcal{O}(\rho^\ell).
        \end{gather}
        \item From \labelcref{FGARCH11CLTCondition}, it follows that $(\sigma_i^2(t);i\in\mathbb{Z})$ and $(y_i(t);i\in \mathbb{Z})$ are weakly dependent processes.
    \end{enumerate}
\end{corollary}
\labelcref{FGARCH11CLTCondition} is known as geometric $L^2$-m-approximibility. $L^p$-m-approximability for functional data was first introduced by \citet{Lpmapproxibility}. A process that is geometrically $L^p$-m-approximable can be approximated by an m-dependent process at a geometric rate. A stochastic process $(X_n)$ is called m-dependent if observations that are more than $m$ time steps apart are independent. More precisely, for any time index $k$, the collection of random variables $\left\{X_n:n\leq k \right\}$ is independent of the collection $\left\{X_n:n\geq k+m \right\}$. In other words, the past up to time $k$ has no influence on the future starting from time $k+m$. Formally, $X_n$ is m-dependent if, for any integer $k$, the $\sigma$-algebra $\mathcal{F}_k^-=\sigma\left\{...,X_{k-2},X_{k-1,X_k} \right\}$ generated by the past observations is independent of the $\sigma$-algebra generated by the future observations $\mathcal{F}_{k+m}^+=\sigma\left\{ X_{k+m},X_{k+m+1},X_{k+m+2},... \right\}$. \citet{Lpmapproxibility} use the concept of $L^p$-m-approximability to
\subsection{Estimation}
\label{subsec:FGARCHEstimation}
\begin{itemize}
    \item Projection onto finite dimensional basis
    \item Least Squares Minimization problem
    \item Solving Least Squares (Connection to uni-/ multivariate GARCH models)
    \item Asymptotic results
    \item (Quasi-Maximum Likelihood Estimation?)
\end{itemize}
%In this next section, the estimation of a functional GARCH(1,1) model will be discussed. To this end, it is assumed that there exists a finite sample of $N$ observed functions $(y_i(t))_{i\in[1,N]}$. In practice, functional objects are typically observed on a fine grid, however, here it is assumed that the functional objects themselves are observed or can be retrieved from the data.\\ My Version before refining with ChatGPT.
In the following section, I address the estimation of a functional GARCH(1,1) model. I consider a setting in which a finite sample of $N$ functional observations $\{y_i(t)\}_{i=1}^N$ is available. Although functional data are commonly recorded on a dense discrete grid in practical applications, I assume here that the underlying functional forms are directly observed or can be reliably reconstructed from the data.\\
Because the objects to be estimated are infinite-dimensional, estimating them directly from a finite sample typically leads to unstable, noisy, and hard-to-interpret results. For this reason, \citet{Aue2017} impose a standard finite-dimensional approximation on the intercept function $\delta(t)$ and on the integral operators $\alpha$ and $\beta$. In particular, they assume that all three objects can be approximated using a fixed number $M$ of orthonormal basis functions $\Phi_M = \{\varphi_1(t), \ldots, \varphi_M(t)\}$ defined on $[0,1]$. Under this assumption, $\delta(t)$ is represented as a finite linear combination of these basis functions, while the operators $\alpha$ and $\beta$ are represented through their kernel functions using the same basis:
\begin{align}
\label{basisdecompositiondelta}
    \delta(t) &= \sum_{m=1}^M d_m \varphi_m(t), \\
    \alpha(t,s) &= \sum_{m,m'=1}^M a_{m,m'} \varphi_m(t)\varphi_{m'}(s), \\
    \label{basisdecompositionbeta}
    \beta(t,s) &= \sum_{m,m'=1}^M b_{m,m'} \varphi_m(t)\varphi_{m'}(s).
\end{align}
From this representation, it is easy to see that estimating the functional objects $\delta(t)$, $\alpha$ and $\beta$ simplifies to estimating $\{d_1,\ldots,d_M,a_{1,1},\ldots,a_{M,M},b_{1,1},\ldots,b_{M,M}\}$. To further simplify notation, \citet{Aue2017} define the following quantities:
\begin{align}
\label{MatricesdAB}
    \mathbf{d}=\begin{pmatrix}
        d_1\\
        \vdots\\
        d_M
    \end{pmatrix}; \quad
    \mathbf{A}=\begin{pmatrix}
        a_{1,1} & \ldots & a_{1,M}\\
        \vdots & \ddots & \vdots\\
        a_{M,1} & \ldots & a_{M,M}
    \end{pmatrix}; \quad
    \mathbf{B}=\begin{pmatrix}
        b_{1,1} & \ldots & b_{1,M}\\
        \vdots & \ddots & \vdots\\
        b_{M,1} & \ldots & b_{M,M}
    \end{pmatrix}
\end{align}
To estimate $d$, $A$, and $B$, it is necessary to project $(y_i(t)^2)_{i\in[1,N]}$ and $(\sigma_i(t))_{i\in[1,N]}$ onto $\Phi_M$ and define the m-dimensional vectors $\mathbf{y}_i^{(2)}=\begin{pmatrix}
    y_{i,1}^{(2)} & \ldots & y_{i,M}^{(2)}
\end{pmatrix}^\top$ and $\mathbf{s}_i^{(2)}=\begin{pmatrix}
    s_{i,1}^{(2)} & \ldots & s_{i,M}^{(2)}
\end{pmatrix}^\top$ where their elements are the result of the inner product with the elements of $\Phi_M$ $y_{i,m}^{(2)}=\langle y_i(t),\varphi_m(t)\rangle$. This leads to the following equation that expresses the projected components of the volatility equation of the FGARCH(1,1) model in matrix form
\begin{align}
\label{FGARCH11MatrixForm}
    s_i^{(2)}=\mathbf{d}+\mathbf{Ay_{i-1}^{(2)}}+\mathbf{Bs_{i-1}^{(2)}}.
\end{align}
Next, note that this equation can be solved recursively to end up with an equation for $ s_i^{(2)}$ that only depends on $y_{i}^{(2)}$
\begin{align}
\label{Matrixrepinftrealparameters}
    s_i^{(2)}=\sum_{\ell=1}^{\infty} \mathbf{B}^{\ell-1}\mathbf{d}+\sum_{\ell=1}^{\infty}\mathbf{B}^{\ell-1}\mathbf{Ay}_{i-\ell}^{(2)}.
\end{align}
Since the true parameters $\theta=(\mathbf{d}^{\top},\mathbf{A}^\top,\mathbf{B}^\top)^\top$ are unknown, \labelcref{Matrixrepinftrealparameters} can be expressed more generally
\begin{align}
    \tilde{s}_i^{(2)}=\tilde{s}_i^{(2)}(\underline{\theta})=\sum_{\ell=1}^{\infty}\underline{\mathbf{B}}^{\ell-1}\underline{\mathbf{d}}+\sum_{\ell=1}^{\infty}\underline{\mathbf{B}}^{\ell-1}\underline{\mathbf{A}}\mathbf{y}_{i-\ell}^{(2)}.
\end{align}
Here, $\underline{\theta}=(\underline{\mathbf{d}}^{\top},\underline{\mathbf{A}}^\top,\underline{\mathbf{B}}^\top)^\top$ is an element of the parameter space $\mathbf{\Theta}\subset\mathbb{R}^{M+2M^2}$. Since in practice only a finite time frame of the series can be observed, $s_i^{(2)}$ is approximated by $\hat{s}_i^{(2)}$
\begin{align}
    \hat{s}_i^{(2)}=\hat{s}_i^{(2)}(\underline{\theta})=\sum_{\ell=1}^{i-1}\underline{\mathbf{B}}^{\ell-1}\underline{\mathbf{d}}+\sum_{\ell=1}^{i-1}\underline{\mathbf{B}}^{\ell-1}\underline{\mathbf{A}}\mathbf{y}_{i-\ell}^{(2)}.
\end{align}
The goal of any objective function, thus, has to be to find $\underline{\hat{\theta}}_n$ such that $\hat{s}_i^{(2)}$ is close to $s_i^{(2)}$. Since $s_i^{(2)}$ is unobservable, \citet{Aue2017} impose
\begin{gather}
\label{UnitVarianceepsilonfunction}
    \mathbb{E}[\varepsilon_i^2(t)]=1; \quad \forall t\in[0,1].
\end{gather}
This is without loss of generality since this condition can always be fulfilled by an appropriate way of scaling. Under this condition, $\mathbb{E}[s_i^{(2)}|\mathcal{F}_{i-1}]=y_i^{(2)}$ where $\mathcal{F}_{i-1}$ is the sigma-algebra generated by all $\varepsilon_\ell:\ell<i$, $\mathcal{F}_{i-1}=\sigma(\varepsilon_\ell:\ell<i)$. As a consequence, \citet{Aue2017} define the following least squares estimator
\begin{align}
\label{ObjectiveFunctionLSFGARCH}
    \underline{\hat{\theta}}_n=argmin\{S_n(\underline{\theta})=\sum_{i=2}^n\{y_i^{(2)}-\hat{s}_i^{(2)}(\underline{\theta}) \}^{\top}\{y_i^{(2)}-\hat{s}_i^{(2)}(\underline{\theta}) \}:\underline{\theta}\in\mathbf{\Theta}\}.
\end{align}
In the following, I will discuss the strong consistency and asymptotic distribution of the least-squares estimator $ \underline{\hat{\theta}}_n$.\\
To prove the strong consistency and asymptotic normality of the estimator, the following Assumption is needed.
\begin{assumption}
\label{Assumption1}
    It is assumed that
    \begin{enumerate}[label=(A\arabic*)]
        \item $\mathbf{y}_1^{(2)}$ is not measurable with respect to $\mathcal{F}_0$,
        \item $\mathbf{\Theta}$ is a compact set and $\theta$ is in the interior of $\mathbf{\Theta}$ and
        \label{Ainvertibility}
        \item there are $0<c_1$ and $c_2<1$ such that $c_1\leq \left|\det(\underline{\mathbf{A}})\right|$ and $\left\| \mathbf{B} \right\|<1$ for all $\underline{\theta}=(\underline{\mathbf{d}}^\top,\underline{\mathbf{A}}^\top,\underline{\mathbf{B}}^\top)\in\mathbf{\Theta}$.
    \end{enumerate}
\end{assumption}
The first part of the assumption rules out trivial predictability of $y_1^{(2)}$ from its past. The second and third parts impose standard regularity conditions on the parameter space. \labelcref{Ainvertibility} also implies that $A$ is non-singular and thus invertible.\\
Using this assumption, the strong consistency of $\underline{\hat{\theta}}_n$ is established in the following Theorem.
\begin{theorem}
\label{StrongConstistencyFGARCH11}
    Assume that $(y_i(t);i\in\mathbb{Z})$ follows a Functional GARCH(1,1) process as defined in \labelcref{FunctionalGARCH11} and \labelcref{FunctionalGARCH11volatility}), that $\delta(t)$ is a non-negative function, and $\alpha$ and $\beta$ are integral operators that map non-negative functions to non-negative functions, that \labelcref{FGARCH11StationarityMomentsCondition} holds with $\nu=1$, that $\mathbb{E}\left[\left\|\varepsilon_0\right\|_2^2\right]<\infty$, and that \labelcref{UnitVarianceepsilonfunction} holds. Then under \Cref{Assumption1} $\underline{\hat{\theta}}_n$ is strongly consistent, meaning that
    \begin{gather}
        \underline{\hat{\theta}}_n \rightarrow \theta
    \end{gather}
    as $n\rightarrow\infty$.
\end{theorem}
To examine the asymptotic distribution of $\underline{\hat{\theta}}_n$, the following quantities are useful
\begin{gather}
    \mathbf{H_0}=\mathbb{E}\left[\frac{\partial \tilde{s}_0^{(2)}(\theta)}{\partial\underline{\theta}}\right], \quad \quad \mathbf{J_0}=\mathbb{E}\left[\left\{y_0^{(2)}-s_0^{(2)}\right\}\left\{y_0^{(2)}-s_0^{(2)}\right\}^\top\right],\\
    \mathbf{Q_0}=\mathbb{E}\left[\left(\frac{\partial \tilde{s}_0^{(2)}(\theta)}{\partial\underline{\theta}}\right)^\top\left(\frac{\partial \tilde{s}_0^{(2)}(\theta)}{\partial\underline{\theta}}\right)\right].
\end{gather}
Under the conditions stated in \Cref{StrongConstistencyFGARCH11}, the matrices $H_0$ and $Q_0$ are well defined. Moreover, if $\mathbb{E}[\left\|\sigma_0^2\right\|_2^4]<\infty$ and $\mathbb{E}[\left\|\varepsilon_0^2\right\|_2^4]<\infty$, the matrix $J_0$ exists. We additionally assume that $Q_0$ is non-singular. The matrix $Q_0$ can be interpreted as the second-moment matrix of the score derivative and, as shown in the Appendix, it is equal to the second derivative of the sample version of $J_0$, the objective function to be minimized. Non-singularity, therefore, guaranties invertibility and ensures that the objective function is locally well behaved.
\begin{theorem}
\label{AsymptoticDistroFGARCH11estimator}
Assume that the conditions of \Cref{StrongConstistencyFGARCH11} are satisfied, that $\mathbf{Q_0}$ is non-singular, and that $\mathbb{E}\!\left[\|y_0^2\|_2^4\right]<\infty$. Then, as $n\to\infty$,
\begin{gather}
\sqrt{n}\,(\underline{\hat{\theta}}_n-\theta)
\;\overset{\mathcal{D}}{\longrightarrow}\;
\mathcal{N}\!\left(0,\,
\mathbf{Q_0}^{-1}\mathbf{H_0}^\top\mathbf{J_0H_0Q_0}^{-1}\right),
\end{gather}
where the limiting distribution is a $(2M^2+M)$-dimensional normal distribution.
\end{theorem}
%The proof of \ref{AsymptoticDistroFGARCH(1,1)estimator} uses the (componentwise) mean value theorem. First, define $\hat{Z}_n(\underline{\theta})=\frac{1}{n}\sum_{i=2}^n\left\{ y_i^{(2)}-\hat{s}_i^{(2)}(\underline{\theta})\right\}^\top \left\{ y_i^{(2)}-\hat{s}_i^{(2)}(\underline{\theta})\right\}$ as the sample version of the least squares value with any $\underline{\theta}$. Denoting the gradient and Hessian by $\frac{\partial \hat{Z}_n(\underline{\theta})}{\partial\underline{\theta}}$ and $\frac{\partial^2 \hat{Z}_n(\underline{\theta})}{\partial\underline{\theta}^2}$, the mean value theorem gives $\frac{\partial \hat{Z}_n(\underline{\theta})}{\partial\underline{\theta}}-\frac{\partial \hat{Z}_n(\theta)}{\partial\underline{\theta}}=\frac{\partial^2 \hat{Z}_n(\check{\theta)}}{\partial\underline{\theta}^2}(\underline{\theta}_n-\theta)$ for some $\check{\theta}$ between $\theta$ and $\underline{\theta}$. Since $\frac{\partial \hat{Z}_n(\underline{\theta})}{\partial\underline{\theta}}=0$, the asymptotic distribution of the estimate $\underline{\theta_n}$ depends only on $\frac{\partial \hat{Z}_n(\theta)}{\partial\underline{\theta}}$ and $\frac{\partial^2 \hat{Z}_n(\check{\theta)}}{\partial\underline{\theta}^2}$. Applying the uniform law of large numbers to $\frac{\partial^2 \hat{Z}_n(\check{\theta)}}{\partial\underline{\theta}^2}$ and the CLT for orthogonal martingale differences to $\frac{\partial \hat{Z}_n(\theta)}{\partial\underline{\theta}}$ yields the asymptotic distribution in \ref{AsymptoticDistroFGARCH(1,1)estimator}. The full proof can be found in the Appendix.\\ My version

The proof of \Cref{AsymptoticDistroFGARCH11estimator} relies on a componentwise application of the mean value theorem. Define
\[
\hat{Z}_n(\underline{\theta})
=
\frac{1}{n}\sum_{i=2}^n
\left\{ y_i^{(2)}-\hat{s}_i^{(2)}(\underline{\theta})\right\}^\top
\left\{ y_i^{(2)}-\hat{s}_i^{(2)}(\underline{\theta})\right\},
\]
the sample least squares criterion for a generic parameter vector $\underline{\theta}$.
Denote its gradient and Hessian by
$\frac{\partial \hat{Z}_n(\underline{\theta})}{\partial\underline{\theta}}$
and
$\frac{\partial^2 \hat{Z}_n(\underline{\theta})}{\partial\underline{\theta}^2}$,
respectively.
Applying the mean value theorem to the score function at the estimator
$\underline{\hat{\theta}}_n$ yields
\[
\frac{\partial \hat{Z}_n(\underline{\hat{\theta}}_n)}{\partial\underline{\theta}}
-
\frac{\partial \hat{Z}_n(\theta)}{\partial\underline{\theta}}
=
\frac{\partial^2 \hat{Z}_n(\check{\theta})}{\partial\underline{\theta}^2}
(\underline{\hat{\theta}}_n-\theta),
\]
for some intermediate value $\check{\theta}$ lying between $\theta$ and
$\underline{\hat{\theta}}_n$.
Since $\frac{\partial \hat{Z}_n(\underline{\theta})}{\partial\underline{\theta}}=0$, the first term vanishes, and thus
\[
\sqrt{n}(\underline{\hat{\theta}}_n-\theta)
=
-
\left(\frac{\partial^2 \hat{Z}_n(\check{\theta})}{\partial\underline{\theta}^2}\right)^{-1}
\sqrt{n}\,
\frac{\partial \hat{Z}_n(\theta)}{\partial\underline{\theta}}.
\]

The asymptotic distribution of $\underline{\hat{\theta}}_n$ therefore depends on the limiting behavior of the score and the Hessian. A uniform law of large numbers yields convergence of the Hessian to $2\mathbf{Q_0}$, while a central limit theorem for orthogonal martingale differences applies to the score evaluated at $\theta$. Combining these results establishes the asymptotic normality stated in \Cref{AsymptoticDistroFGARCH11estimator}. The full proof is deferred to the Appendix.\\
\citet{Cerovecki2019} derive an alternative estimate based on (pseudo) quasi-maximum likelihood. Applying maximum likelihood techniques to functional data is in general not straightforward since, in general, there is no density function for curves. As a consequence, \citeauthor{Cerovecki2019} also assume that $\delta(t)$, $\alpha$, and $\beta$ can be decomposed into a finite combination of baselines functions similar to (\ref{basisdecompositiondelta}-\ref{basisdecompositionbeta}). Notably, they do not assume that these baseline functions are orthonormal, but rather linear independent and non-negative. They show that under these conditions they arrive at a similar as \labelcref{FGARCH11MatrixForm} which sequentially reduces to a semi-strong multivariate CCC-GARCH model with $R=\mathbb{I}_M$, where $\mathbb{I}_M$ is the identity matrix with dimensions $M\times M$.\\
They prove the consistency and asymptotic normality of the quasi-maximum likelihood estimator for this semi-strong multivariate CCC-GARCH model. Following this result, they prove the consistency and asymptotic normality of the (pseudo) quasi-maximum likelihood estimator of the Functional GARCH model. The estimator for the Functional GARCH model is then defined as follows.
\begin{gather}
    \underline{\hat{\theta}}_n= \arg \min_{\theta \in \Theta} \sum_{i=1}^n \tilde{\ell}_i(\theta), \quad \tilde{\ell}_i(\theta)=\sum_{m=1}^M\left\{\frac{\langle y_i^2,\varphi_m\rangle}{\langle \tilde{\sigma}_i^2,\varphi_m\rangle}+\log\langle \tilde{\sigma}_i^2,\varphi_m\rangle \right\},
\end{gather}
where the empirical volatility is calculated recursively from\\
Finally, \citeauthor{Aue2017} discuss the estimation of basis functions. In general, any orthonormal basis can be used as basis functions. Popular choices include B-splines, Fourier bases, and wavelet functions. An alternative approach is to estimate Functional Principal Components from the data and use these as basis functions. 
\subsection{Value-at-risk}
\label{subsec:ValueAtRisk}
%\begin{itemize}
%    \item Explanation of the concept of value-at-risk
%\end{itemize}
In this section, I will give a quick overview of the concept of value-at-risk. I will focus on the definition of value-at-risk as well as the criticisms that the concept of value-at-risk has faced.\\
Value-at-risk is intended to specify a limit below which the value of, for example, a security cannot fall with a predetermined probability. In probability terms, this means that under the assumption that the map $x\rightarrow \mathbb{P}(y_i\leq x|\mathcal{F}_{i-1})$ is strictly increasing where $\mathcal{F}_{i-1}$ denotes the $\sigma$-algebra up to $i-1$
\begin{gather}
    \mathbb{P}(y_i\leq \text{VaR}_i^\tau|\mathcal{F}_{i-1})=\tau.
\end{gather}
$\text{VaR}_i^\tau$ here denotes the almost surely unique $\mathcal{F}_{i-1}$-measurable random variable at quantile $\tau$.\\
From this definition, a few properties of value-at-risk can be inferred. 
The concept of value-at-risk has been criticized a lot. As a result, extensions to the definition given above have been established. The main points of critic to value-at-risk are...
\section{Simulation Study}
\label{sec:Simulation}
%\begin{itemize}
%    \item Explain bootstrapping and Probability scoring measures
%\end{itemize}
In the next section, I will explore the finite sample properties of the Functional GARCH model and highlight a few practical considerations that have to be made when using the Functional GARCH model for forecasting intra-day value-at-risk in practice.\\
First, it should be emphasized that functional data are, in practice, observed on a finite grid. Consequently, both the properties discussed in \Cref{subsec:FGARCHModel} and the estimation procedure described in \Cref{subsec:FGARCHEstimation} are implemented in a discretized framework. This discretization does not yield a purely multivariate GARCH model; rather, it leads to a finite-dimensional approximation of the underlying functional model. In particular, operations defined in the functional setting—such as integrals in inner products—must be approximated numerically, which requires explicit consideration of the grid structure and the distances between adjacent grid points. For all simulations, functional data are approximated on an equidistant grid between $[0,1]$ with 200 grid points.\\
Second, to forecast value-at-risk, an appropriate method has to be used to approximate the pointwise distribution of the innovation process $\varepsilon_i(t)$. Then the value-at-risk forecast is equal to $\hat{\sigma}_i(t)\hat{\varepsilon}_i^\tau(t)$ where $\hat{\sigma}_i(t)$ is the forecasted value of $\sigma_i(t)$ using the fitted Functional GARCH model and $\hat{\varepsilon}_i^\tau(t)$ is the estimated $\tau$ quantile of $\varepsilon_i(t)$. To estimate $\hat{\varepsilon}_i^\tau(t)$ \citet{RICE20201023} propose three methods. For the first method, \citeauthor{RICE20201023} impose that $\varepsilon_i(t)$ follows a Ornstein-Uhlenbeck process. The second method estimates $\hat{\varepsilon}_i^\tau(t)$ using a bootstrap algorithm. Method three utilizes extreme value theory to better model the tails of the distribution. I will focus on the second method since it performed the best in practice \citep{RICE20201023}.\\
To obtain $\hat{\varepsilon}_i^\tau(t)$ using the bootstrap algorithm first the fitted $\hat{\sigma}_i(t)$ needs to be computed. Next, $\hat{\varepsilon}_i(t)$ is calculated by $\hat{\varepsilon}_i(t)=\frac{y_i(t)}{\hat{\sigma}_i(t)}$. Finally, to get $\hat{\varepsilon}_i^\tau(t)$ the estimated error terms are resampled $B$ many times with replacement. $\hat{\varepsilon}_i^\tau(t)$ is then equal to the pointwise $\tau$-quantile of $(\hat{\varepsilon}_i^{(b)}(t))_{b\in[1,B]}$.\\
Thirdly, to properly evaluate the results of the functional GARCH model I will employ two statistical tests.
\subsection{Simulation 1}
\label{subsec:Simulation1}
\begin{itemize}
    \item Similar to Aue et al. Show simulation where the method works.
\end{itemize}
\begin{align}
    \varepsilon_i(t)=\frac{\sqrt{\log2}}{2^{200t}}\cdot B_i \left( \frac{2^{400t}}{\log2}\right); (B_i:i\in \mathbb{Z})
\end{align}
\subsection{Simulation 2}
\label{subsec:Simulation2}
\begin{itemize}
    \item Simulation 2 should demonstrate some of the weaknesses from the fGARCH model either by violating one of its assumptions or/and by giving reasons why the fGARCH model should be extended in some type of way (transition to fGARCH-X or fGARCH with multiple constant parameters or similar)
\end{itemize}
\section{Discussion}
\label{sec:Discussion}
\begin{itemize}
    \item Summarize thesis
    \item list possible extensions and literature that has already extended FGARCH
    \item (Give outlook on future of risk measures? (Expected Shortfall?))
\end{itemize}
In this thesis, I have presented the Functional GARCH model and
\newpage
{\setstretch{0.6}
\bibliographystyle{agsm}
\bibliography{bibliography}
}

\appendix
\section{Appendix}
\label{sec:Appendix}
%\renewcommand{\thesubsection}{\Alph{subsection}}
\subsection{Proofs of Section 2.1}
\begin{align}
    \varepsilon_i(t) = e^{-\frac{t}{2}}W_i(e^t)
\end{align}
Let $\sum a_n$ be a series with non-negative terms and let $\lambda = \varlimsup a_n^{1/n}$. Then (i) if $\lambda<1$ the series $\sum a_n$ converges, (ii) if $\lambda>1$ the series $\sum a_n$ diverges.
\subsection{Proofs of Section 2.2}
\begin{align}
    (\Gamma_{i,1}x)(t)&=\int \gamma_{i-1}(t,s)x(s)ds \\
    \lVert\Gamma_{i,1}\rVert_{\mathcal{N}}&\leq \lVert\gamma_{i-1}\rVert_2
\end{align}
\subsection{Proofs of Section 2.3}
\begin{gather}
    \tilde{Z}_n(\underline{\theta})=\frac{1}{n}\sum_{i=2}^n\left\{ y_i^{(2)}-\tilde{s}_i^{(2)}(\underline{\theta})\right\}^\top \left\{ y_i^{(2)}-\tilde{s}_i^{(2)}(\underline{\theta})\right\} \rightarrow \sum_{m=1}^M\mathbb{E}\left[\langle\sigma_0^2(\varepsilon_0^2-1),\varphi_m \rangle^2 \right]+M(\underline{\theta}) \quad \text{a.s.}
\end{gather}
\subsection{Comparison discretized FGARCH(1,1) and PC-GARCH}
Discretizing the FGARCH(1,1) model results in the model simplifying to
\begin{gather}
    y_i=\sigma_i \odot \varepsilon_i\\
    \label{DiscretizedFGARCH11}
    \sigma_i^2 = \delta+\alpha y_{i-1}^2+\beta\sigma_{i-1}^2
\end{gather}
where $\odot$ denotes the Hadamard product, $y_i,\sigma_i,\varepsilon_i,\sigma_i^2,y_i^2$ and $\delta$ are all vectors with length $T$ where $T$ is the number of grid points. $\alpha$ and $\beta$ denote $T \times T$ matrices. $\sigma_i^2$ and $y_i^2$ are the pointwise squared values of $\sigma_i$ and $y_i$ so $\sigma_i^2=\sigma_i\odot \sigma_i$ and $y_i^2=y_i\odot y_i$.\\
With the assumption that $\delta$,$\alpha$ and $\beta$ can all be decomposed into a finite set of basis functions that can be consistently estimated by FPCA, \labelcref{DiscretizedFGARCH11} can further be rewritten into
\begin{gather}
    s_i^{(2)}=\mathbf{d}+\mathbf{Ay_{i-1}^{(2)}}+\mathbf{Bs_{i-1}^{(2)}}
\end{gather}
which is equivalent to \labelcref{FGARCH11MatrixForm}. $s_i^{(2)}$ is then estimated and projected back to $\sigma_i^2$ by
Now I will rewrite the PC-GARCH model to show the similar structure the two models possess. As before, the Principal Components GARCH (PC-GARCH) model is defined by
\begin{gather}
\label{PCGARCHyequationAppendix}
    y_t = H_t^{\frac{1}{2}}\varepsilon_t,\\
    \label{PCGARCHHequationAppendix}
    H_t = P \Lambda_t P^{\top}.
\end{gather}
Here, $P$ is an $m \times m$ orthogonal matrix, while $\Lambda_t$ is an $m \times m$ diagonal matrix with diagonal elements $\text{diag}(\lambda_{1,t}, \lambda_{2,t}, \ldots, \lambda_{m,t})$. The matrix $P$ is typically estimated using Principal Component Analysis, and the diagonal elements of $\Lambda_t$ are assumed to follow univariate GARCH(1,1) processes:
\begin{gather}
\label{PCGARCHlambdaequationAppendix}
    \lambda_{i,t} = \delta_i + \alpha_i f_{i,t-1}^{2} + \beta_i \lambda_{i,t-1}, \quad i = 1, \ldots, m.
\end{gather}
Expanding \labelcref{PCGARCHHequationAppendix} by \labelcref{PCGARCHlambdaequationAppendix} results into
\begin{gather}
    P \Lambda_t P^{\top} = P \mathbf{D} P^{\top}+
\end{gather}
\end{document}
